% ---
% id: reductions/rem:dominating
% kind: rem
% title: "Dominated detectors and device families"
% summary: >-
%   A detector whose POVM dominates a characterized one is the characterized
%   device plus a weighted flag, as in the published block-resolved
%   construction. By [[prop:common]], a family of devices whose members squash
%   to a common target becomes one squashed model with the label a branch
%   register.
% prominence: standard
% status: accepted
% uses: [prop:common]
% mentions: [cor:instances, lem:master]
% assumes: []
% refs:
%   - "nahar2025imperfect, Thm. 3, Eqs. (20)-(23)"
%   - "douglas1966majorization, Thm. 1"
% provenance: newly-derived
% lean:
% ---

\begin{remark}\label{rem:dominating}
A device has the form (i) of Corollary~\ref{cor:instances} whenever its POVM
dominates a characterized one: if $M_v\succeq(1-q)G_v$ for all $v$ with
$q\in(0,1]$, then $R_v\coloneqq\big(M_v-(1-q)G_v\big)/q$ is positive and sums
to $\idon{\Bsys}$, so any detector that is $q$-close to a characterized
detector \emph{in this operator sense} may be analysed as the characterized
device plus a $q$-weighted flag. The published block-resolved form of that
construction is \cite[Thm.~3, Eqs.~(20)--(23)]{nahar2025imperfect}, which
builds exactly this channel from the block dominations
$\Gamma_{m,\text{noise}}-(1-q_m)\Gamma_{m,\text{ideal}}\succeq0$ and records
the weight $\Lambda^{\dagger}[\Pi_m]=(1-q_m)\Pi_m$ kept on block $m$, the flag
receiving $q_m$.

By Proposition~\ref{prop:common}, a family of devices indexed by a label $\xi$
that the adversary selects becomes one squashed model as soon as each member
squashes to a \emph{common} target, with the label a branch register carrying
its own constraints from Lemma~\ref{lem:master}\,(iv). The operator
factorization such a construction rests on (range inclusion, majorization and
factorization are equivalent for bounded operators) is
\cite[Thm.~1]{douglas1966majorization}; no proof of this contract uses it.
\end{remark}
